\documentclass[12pt,a4paper]{article}

\usepackage{cmap}
\usepackage[T2A]{fontenc}
\usepackage[utf8]{inputenc}
\usepackage[russian]{babel}
\usepackage{amsmath,amssymb}
\usepackage[a4paper,margin=2cm]{geometry}

\newcommand{\R}{\mathbb{R}}
\newcommand{\eps}{\varepsilon}
\newcommand{\Ball}[2]{U_{#1}(#2)}
\newcommand{\PBall}[2]{\mathring U_{#1}(#2)}

\setlength{\parindent}{0pt}
\setlength{\parskip}{0.45em}
\sloppy

\begin{document}

\begin{center}
{\Large\bfseries Билет 1}
\end{center}

\noindent\fbox{%
  \begin{minipage}{\dimexpr\linewidth-2\fboxsep-2\fboxrule\relax}
  \normalfont
Предел последовательности точек в $\mathbb{R}^n$. Связь между сходимостью
последовательности точек и сходимостью последовательностей их координат.
Внутренние, предельные, изолированные точки множества. Открытые и замкнутые
множества, их свойства. Непрерывная кривая, область, компакт. Критерий
компактности в $\mathbb{R}^n$. Внутренность, замыкание и граница множества.
  \end{minipage}%
}

\section*{Краткий план}

Работаем в евклидовом пространстве \(\R^n\) с нормой
\[
  |x|=\sqrt{x_1^2+\dots+x_n^2}, \qquad \rho(x,y)=|x-y|.
\]
Сначала вводятся окрестности и предел последовательности точек. Затем
доказывается координатный критерий сходимости. После этого вводятся
внутренние, предельные, изолированные и граничные точки, открытые и
замкнутые множества, кривая и область. В конце формулируется и доказывается
критерий компактности: в \(\R^n\) компактность эквивалентна замкнутости и
ограниченности.

\section*{Определения}

\textbf{Окрестность.}
Для \(x_0=(x_0^1,\dots,x_0^n)\in\R^n\)
\[
  \Ball{\eps}{x_0}=\{x\in\R^n:\ |x-x_0|<\eps\}.
\]
Проколотая окрестность:
\[
  \PBall{\eps}{x_0}=\Ball{\eps}{x_0}\setminus\{x_0\}.
\]

\textbf{Предел последовательности точек.}
Последовательность \(x_k\in\R^n\) сходится к \(x_0\in\R^n\), если
\[
  \forall \eps>0\ \exists N\ \forall k\ge N:\ |x_k-x_0|<\eps.
\]
Пишут \(x_k\to x_0\) или \(\lim\limits_{k\to\infty}x_k=x_0\).

\textbf{Внутренняя точка и внутренность.}
Точка \(x_0\in\R^n\) называется внутренней точкой множества \(X\subset\R^n\),
если
\[
  \exists \eps>0:\ \Ball{\eps}{x_0}\subset X.
\]
Множество всех внутренних точек называется внутренностью и обозначается
\(\operatorname{int}X\).

\textbf{Открытое множество.}
Множество \(X\) называется открытым, если каждая его точка внутренняя:
\[
  X=\operatorname{int}X.
\]
Пустое множество считается открытым.

\textbf{Точка прикосновения и замыкание.}
Точка \(x_0\in\R^n\) называется точкой прикосновения \(X\), если
\[
  \forall \eps>0:\ \Ball{\eps}{x_0}\cap X\ne\varnothing.
\]
Множество всех точек прикосновения называется замыканием \(X\) и обозначается
\(\overline X\).

\textbf{Замкнутое множество.}
Множество \(X\) называется замкнутым, если оно содержит все свои точки
прикосновения:
\[
  \overline X\subset X,
\]
то есть, с учетом \(X\subset\overline X\), если \(\overline X=X\). Пустое
множество считается замкнутым.

\textbf{Предельная и изолированная точки.}
Точка \(x_0\in\R^n\) называется предельной точкой множества \(X\), если
\[
  \forall \eps>0:\ \PBall{\eps}{x_0}\cap X\ne\varnothing.
\]
Эквивалентно: существует последовательность \(x_k\in X\), \(x_k\ne x_0\),
такая что \(x_k\to x_0\).

Точка \(x_0\) называется изолированной точкой \(X\), если
\[
  x_0\in X
  \quad\text{и}\quad
  \exists \eps>0:\ \PBall{\eps}{x_0}\cap X=\varnothing.
\]
Точка прикосновения является либо предельной точкой, либо изолированной
точкой самого множества.

\textbf{Граница.}
Границей множества \(X\) называется
\[
  \partial X=\overline X\setminus \operatorname{int}X.
\]
Эквивалентно,
\[
  x_0\in\partial X
  \Longleftrightarrow
  \forall \eps>0:\ \Ball{\eps}{x_0}\cap X\ne\varnothing
  \ \text{и}\ 
  \Ball{\eps}{x_0}\setminus X\ne\varnothing.
\]

\textbf{Непрерывная кривая.}
Кривой в \(\R^n\) называется годограф непрерывной вектор-функции
\[
  r:[a,b]\to\R^n,\qquad
  \Gamma=\{r(t):\ t\in[a,b]\}.
\]

\textbf{Область.}
\emph{Стандартное уточнение.} Областью в \(\R^n\) называют открытое линейно
связное множество. Линейная связность означает, что любые две точки
\(x_1,x_2\in X\) можно соединить кривой, лежащей в \(X\): существует
непрерывная \(r:[t_1,t_2]\to X\), такая что \(r(t_1)=x_1\), \(r(t_2)=x_2\).

\textbf{Ограниченность.}
Множество \(X\subset\R^n\) ограничено, если
\[
  \exists C\in\R\ \forall x\in X:\ |x|\le C.
\]

\textbf{Компакт.}
В программе используется последовательностное определение: множество
\(X\subset\R^n\) называется компактом, если из любой последовательности
\(x_k\in X\) можно выделить подпоследовательность, сходящуюся к некоторому
элементу \(X\).

\section*{Теоремы и свойства}

\textbf{1. Координатный критерий сходимости.}
Пусть
\[
  x_k=(x_k^1,\dots,x_k^n), \qquad x_0=(x_0^1,\dots,x_0^n).
\]
Тогда
\[
  x_k\to x_0
  \Longleftrightarrow
  \forall i=1,\dots,n:\ x_k^i\to x_0^i.
\]

\textbf{2. Критерий точки прикосновения.}
Для \(X\subset\R^n\)
\[
  x_0\in\overline X
  \Longleftrightarrow
  \exists\{x_k\}\subset X:\ x_k\to x_0.
\]

\textbf{3. Основные свойства внутренности и замыкания.}
Для любого \(X\subset\R^n\)
\[
  \operatorname{int}X\subset X\subset\overline X.
\]
Кроме того:
\[
  X\text{ открыто}\Longleftrightarrow X=\operatorname{int}X,\qquad
  X\text{ замкнуто}\Longleftrightarrow X=\overline X.
\]
Внутренность \(\operatorname{int}X\) открыта, а замыкание \(\overline X\)
замкнуто.

\textbf{4. Дополнение замкнутого множества.}
\[
  X\text{ замкнуто}\Longleftrightarrow \R^n\setminus X\text{ открыто}.
\]
Отсюда следуют стандартные свойства: любое объединение открытых множеств
открыто; конечное пересечение открытых множеств открыто; любое пересечение
замкнутых множеств замкнуто; конечное объединение замкнутых множеств
замкнуто.

\textbf{5. Теорема Больцано--Вейерштрасса в \(\R^n\).}
Из любой ограниченной последовательности в \(\R^n\) можно выделить сходящуюся
подпоследовательность.

\textbf{6. Критерий компактности в \(\R^n\).}
Множество \(X\subset\R^n\) является компактом тогда и только тогда, когда
\(X\) ограничено и замкнуто.

\section*{Доказательства}

\textbf{Координатный критерий.}
Пусть \(x_k\to x_0\). Тогда для каждой координаты
\[
  |x_k^i-x_0^i|\le |x_k-x_0|\to 0,
\]
значит \(x_k^i\to x_0^i\).

Обратно, пусть \(x_k^i\to x_0^i\) для всех \(i=1,\dots,n\). Тогда
\[
  |x_k-x_0|^2
  =\sum_{i=1}^n (x_k^i-x_0^i)^2 \to 0,
\]
поскольку конечная сумма бесконечно малых стремится к нулю. Следовательно,
\(|x_k-x_0|\to0\), то есть \(x_k\to x_0\).

\textbf{Критерий точки прикосновения.}
Если \(x_k\in X\) и \(x_k\to x_0\), то всякая окрестность \(x_0\) содержит
некоторый \(x_k\), значит пересекает \(X\), и \(x_0\in\overline X\).

Обратно, если \(x_0\in\overline X\), то для каждого \(k\in\mathbb N\) можно
выбрать
\[
  x_k\in X\cap\Ball{1/k}{x_0}.
\]
Тогда \(|x_k-x_0|<1/k\to0\), следовательно, \(x_k\to x_0\).

\textbf{Открытость шаров.}
Пусть \(x\in\Ball{\eps}{x_0}\). Положим
\[
  \delta=\eps-|x-x_0|>0.
\]
Если \(y\in\Ball{\delta}{x}\), то по неравенству треугольника
\[
  |y-x_0|\le |y-x|+|x-x_0|<\delta+|x-x_0|=\eps.
\]
Значит \(\Ball{\delta}{x}\subset\Ball{\eps}{x_0}\), поэтому каждая точка шара
внутренняя, и шар открыт.

\textbf{Свойства внутренности и замыкания.}
Включение \(\operatorname{int}X\subset X\) следует из определения: если
\(\Ball{\eps}{x}\subset X\), то \(x\in X\). Включение \(X\subset\overline X\)
следует из того, что любая окрестность точки \(x\in X\) пересекает \(X\) хотя
бы в точке \(x\).

Покажем, что \(\operatorname{int}X\) открыта. Если
\(x\in\operatorname{int}X\), то для некоторого \(\eps>0\)
\(\Ball{\eps}{x}\subset X\). Так как шар открыт, каждая его точка имеет малую
окрестность, лежащую в \(\Ball{\eps}{x}\), а значит и в \(X\). Поэтому
\(\Ball{\eps}{x}\subset \operatorname{int}X\), и \(x\) внутренняя точка
\(\operatorname{int}X\).

Покажем, что \(\overline X\) замкнуто. Пусть \(x\) является точкой
прикосновения \(\overline X\). Для любого \(\eps>0\) найдется
\(y\in\overline X\cap \Ball{\eps/2}{x}\). Так как \(y\in\overline X\), найдется
\(z\in X\cap\Ball{\eps/2}{y}\). Тогда
\[
  |z-x|\le |z-y|+|y-x|<\eps.
\]
Значит любая \(\eps\)-окрестность \(x\) пересекает \(X\), то есть
\(x\in\overline X\). Следовательно, \(\overline X\) замкнуто.

\textbf{Критерий границы.}
По определению \(x\in\partial X\) означает
\[
  x\in\overline X
  \quad\text{и}\quad
  x\notin\operatorname{int}X.
\]
Первое условие значит, что каждая окрестность \(x\) пересекает \(X\). Второе
значит, что никакая окрестность \(x\) не лежит целиком в \(X\), то есть в
каждой окрестности есть точка из дополнения \(\R^n\setminus X\). Это и дает
эквивалентное описание граничной точки.

\textbf{Замкнутость и открытость дополнения.}
Если \(X\) замкнуто и \(x\notin X\), то \(x\notin\overline X\). Значит для
некоторого \(\eps>0\)
\[
  \Ball{\eps}{x}\cap X=\varnothing,
\]
то есть \(\Ball{\eps}{x}\subset\R^n\setminus X\). Поэтому дополнение открыто.

Обратно, если \(\R^n\setminus X\) открыто и \(x\in\overline X\), то \(x\notin
\R^n\setminus X\), иначе некоторая окрестность \(x\) лежала бы в дополнении и
не пересекала бы \(X\). Значит \(x\in X\), поэтому \(\overline X\subset X\),
и \(X\) замкнуто.

\textbf{Теорема Больцано--Вейерштрасса в \(\R^n\).}
Доказательство проводится индукцией по \(n\). При \(n=1\) это классическая
теорема для числовых последовательностей. Пусть утверждение верно в
\(\R^m\). Рассмотрим ограниченную последовательность в \(\R^{m+1}\):
\[
  x_k=(x_k^1,\dots,x_k^m,x_k^{m+1}).
\]
Первые \(m\) координат образуют ограниченную последовательность в \(\R^m\),
по предположению индукции из нее можно выделить подпоследовательность, у
которой первые \(m\) координат сходятся. Последняя координата этой
подпоследовательности является ограниченной числовой последовательностью,
поэтому из нее можно выделить сходящуюся подпоследовательность. У полученной
подпоследовательности сходятся все координаты, значит по координатному
критерию она сходится в \(\R^{m+1}\).

\textbf{Критерий компактности.}
Пусть \(X\) ограничено и замкнуто. Возьмем любую последовательность
\(x_k\in X\). Она ограничена, поэтому по теореме Больцано--Вейерштрасса из
нее можно выделить подпоследовательность \(x_{k_j}\to x_0\in\R^n\). Так как
\(x_{k_j}\in X\), точка \(x_0\) является точкой прикосновения \(X\), то есть
\(x_0\in\overline X\). Множество \(X\) замкнуто, значит
\(\overline X=X\), поэтому \(x_0\in X\). Следовательно, \(X\) компактно.

Обратно, пусть \(X\) компактно. Если \(X\) неограничено, то для каждого
\(k\in\mathbb N\) можно выбрать \(x_k\in X\) так, что \(|x_k|>k\). У такой
последовательности нет сходящейся подпоследовательности, потому что всякая
сходящаяся последовательность ограничена. Это противоречит компактности.
Значит \(X\) ограничено.

Если \(X\) не замкнуто, то существует \(x_0\in\overline X\setminus X\). По
критерию точки прикосновения существует последовательность \(x_k\in X\),
такая что \(x_k\to x_0\). Любая ее подпоследовательность тоже сходится к
\(x_0\), но \(x_0\notin X\). Значит нельзя выделить подпоследовательность,
сходящуюся к элементу \(X\), что противоречит компактности. Поэтому \(X\)
замкнуто.

\section*{Финальная устная версия}

В \(\R^n\) сходимость определяется через евклидову метрику:
\[
  x_k\to x_0 \Longleftrightarrow
  \forall\eps>0\ \exists N\ \forall k\ge N:\ |x_k-x_0|<\eps.
\]
Главный критерий: последовательность точек сходится тогда и только тогда,
когда сходятся все координатные последовательности. Это следует из оценок
\(|x_k^i-x_0^i|\le |x_k-x_0|\) и формулы
\(|x_k-x_0|^2=\sum_i(x_k^i-x_0^i)^2\).

Внутренняя точка множества имеет целую окрестность внутри множества;
внутренность \(\operatorname{int}X\) есть множество всех таких точек. Точка
прикосновения -- точка, каждая окрестность которой пересекает \(X\);
замыкание \(\overline X\) есть множество всех точек прикосновения.
Предельная точка требует пересечения с \(X\) в каждой проколотой
окрестности, а изолированная точка принадлежит \(X\), но имеет проколотую
окрестность без точек \(X\). Граница задается формулой
\(\partial X=\overline X\setminus\operatorname{int}X\), то есть в любой
окрестности граничной точки есть и точки \(X\), и точки дополнения.

Множество открыто, если \(X=\operatorname{int}X\), и замкнуто, если
\(X=\overline X\). Внутренность всегда открыта, замыкание всегда замкнуто,
а \(X\) замкнуто тогда и только тогда, когда \(\R^n\setminus X\) открыто.

Кривая -- образ непрерывной функции \(r:[a,b]\to\R^n\). Область -- открытое
линейно связное множество, то есть любые две ее точки соединяются кривой,
лежащей в этой области.

Компакт в \(\R^n\) -- множество, из любой последовательности точек которого
можно выделить подпоследовательность, сходящуюся к точке этого же множества.
Критерий компактности: \(X\subset\R^n\) компактно тогда и только тогда,
когда \(X\) ограничено и замкнуто. Доказательство: из ограниченности по
Больцано--Вейерштрассу выделяем сходящуюся подпоследовательность, а
замкнутость гарантирует, что ее предел лежит в \(X\). Обратно, компакт не
может быть неограниченным, иначе выбираем \(|x_k|>k\), и не может быть
незамкнутым, иначе последовательность из \(X\) сходится к точке
\(\overline X\setminus X\), что противоречит определению компактности.

\section*{Источники}

Иванов Г.Е. \emph{Лекции по математическому анализу. Часть 1}:
\texttt{IvGE\_1.pdf}, стр. 144--145, 158--164, 173, 177, 335;
\texttt{IvGE\_2.pdf}, стр. 56--57.

\end{document}
